\chapter[Sermon 68. The Seven Angels Are Described, Coming Forth to Execute the Seven Plagues.]{The Seven Angels Are Described, Coming Forth to Execute the Seven Plagues.}
\label{sermon:068}
\markright{Sermon 68. The Seven Angels Are Described, Coming Forth to ...}

And after that I looked, and behold, the Temple of the Tabernacle of testimony was open in Heaven, and the seven Angels came out of the Temple, which had the seven plagues, clothed in pure and bright linen, and having their breasts girded with golden girdles. And one of the four beasts gave unto the seven Angels seven golden vials full of the wrath of God, who lives forevermore. And the temple was full of smoke, from the glory of God, and from his power: and no man was able to enter into the Temple, until the seven plagues of the seven Angels were fulfilled.

\index{Repentance}He now returns to the description of God’s judgments, from which he had made a brief digression. This treatise is very fruitful. For God’s judgments are the punishments or pains of the wicked, and testimonies of God’s righteousness and truth. Again, the godly are strengthened in their hope. For they see that not one jot falls from the words and threatenings of God, even though he may be long-suffering, tolerates them for a time, and even seems to favor and spare the ungodly. Therefore, the godly perceive that their hope is not vain. They also learn to fear God and to pray continually, lest, being intoxicated by the pleasures and felicities of this world, they turn away from God to ungodliness. Finally, the wicked are terrified by these pains, are provoked to repentance, which, if they refuse, they will undoubtedly feel plagues, as Pharaoh did.

\index{Hebrews, Epistle to the}But before the Anglo-Saxons’ power brings forth the cups of plagues received, they are most fiercely and diligently described. And it is shown from whence they came, that is, what is the origin of God’s judgments. They come out of the temple set open, and that out of the temple of the Tabernacle of Witness, which is in heaven. For Moses saw a temple on the mount, and that also in heaven, after the similitude whereof he was commanded of God to make the tabernacle of Witness. Therefore, the tabernacle of Witness was fashioned and built after the shape exhibited and seen in heaven, which the blessed Apostle to the Hebrews calls the pattern, that is, the very example or patron. For it was said to Moses, “See that you make everything according to the pattern which was shown you on the Mount.” Which thing Moses did accordingly. But such things as came forth of the Tabernacle of Witness made on earth, seemed to the Israelites just and holy. Hereof were delivered the oracles and answers of God, which it was not lawful to speak against. Therefore, when we hear now that the very judgments of God against the wicked world, pains and punishments, come out of the true temple itself, the pattern I mean and that celestial, who should hereafter doubt that all the judgments of God, with which he plagues the ungodly, be sacred and holy? And whilst the ungodly are plagued, that we must think nothing else, but that a sentence, as it were an oracle, is come or pronounced from heaven, which it is unlawful to gainsay? To conclude, the divine judgments do proceed out of the throne of God, wherefore they can not but be most holy. Otherwise, we shall hear in the twenty-first chapter that there is no temple in heaven. These are therefore types and figures, not matters true and permanent; but after they have signified this, for which they were instituted, they are passing and fading away.

\index{Resurrection}Hereunto also appertains the apparel of Angels, that hereof men may also esteem the judgments of God. They are said to be clothed in pure linen, or clean and white, or bright, whereby is signified that the judgments of God are unspotted and bright. For we have heard that these things which John saw were signs. Therefore we may not imagine carnal things in heavenly matters, but spiritually to expound such things as in the sign seem to be as it were corporal. The garment in this world is changed with the state of things. For they use white garments in victories and triumphs, black at burials and mournings, red in battle. Here is signified therefore that the judgments of God are most pure, and that God overcomes and triumphs over the ungodly. At the resurrection and ascension of our Lord Angels appeared in white garments, and shining bright, to signify the glory of Christ. Now is the very breast girded with a girdle, and that in deed with a golden girdle. Gold is a token of purity. In the breast is the seat of affections. The girdle binds, moreover prepares for the journey. Therefore it betokens that the judgments of God are prepared, and in a readiness: the same to restrain affections, that is to say, not to be pronounced or done of envy or malice, love or favour, but to be just, moderate and upright.

And one of the beasts gave to the seven angels, revengers and punishers, seven bowls, and these were full of God’s wrath. Now, although God does not need the help of creatures, nor receives anything from them as if lacking anything, yet he did not make his creatures in vain, and he acts in order. All creatures, without doubt (for I said in chapter four how the beasts signify the universality of creatures), contribute their labor against the wicked, and whatever they have from God (and they have all things) at his will and commandment, they willingly and freely employ to execute God’s judgments. Thus, fire falling from heaven upon Sodom and the cities around it, ministered the plague or cup of God’s wrath to the angel revenger. So the water overwhelmed Pharaoh and his host. So the earth opened and swallowed the company of Dathan and Abiram, and so forth. Thus, the armies of the Gentiles employ themselves to punish the ungodly. The walls of Jericho fall, the hail destroys the Canaanites. Thus God, without any difficulty, punishes his enemies, saying that all creatures are ready to aid and assist. And the bowl or cup is of gold, for again is signified the justice and equity of God’s judgments. And where God is called a ruler living forever, his eternity and majesty are signified, which neither the transitory things of this world nor human infirmities shall overcome. In the sight of the living God, all the wicked shall fall and perish eternally.

\index{Primasius}After this, the Apostle sees the temple filled with smoke for the majesty of God and for his power. That smoke is a sign of God’s presence, it appears by many places in Scripture, but chiefly in chapter 8 of the 3rd book of Kings. It is also a token of God’s wrath. For Aretha says, “smoke is a token of God’s wrath,” according as it is said, “smoke ascended in his wrath.” And neither is smoke without fire, nor fire without smoke. Moreover, smoke hurts the eyes and makes them blind. So in Isaiah 6, the temple of God, which Isaiah sees, is filled with smoke. And at this present, not only the presence of God and of his wrath appear to be signified, but also to be figured, that the judgments of God are unsearchable, so that the things which he himself reveals not to us, we cannot attain to. For his majesty is infinite, and his power passes all things. Primasius, Bishop of Utica in Africa, expounding this place, says, “Think that smoke signifies that all men cannot penetrate the secrets of God’s judgments, and that the eyes and minds of mortal men shall, at the contemplation of the plagues inflicted, grope in darkness.” Which now he determines to utter, and unto the final end of the same, he affirms the smoke to abide still in the temple. Thus says he.

It now appears to explain the same, and no man could enter into the Temple, and so forth. But it is certainly true, by the authority of the evangelical and apostolic doctrine, that the souls departing from the body before the end and final judgment go directly into the blessed places and there experience the joys promised by God, which is true. Therefore, another thing is signified, namely, that before the end of all things, the saints cannot clearly see all of God’s judgments. For here we see through a glass, but there face to face, and we shall know God himself, and the truth and manner of his judgments. Primasius, neither could any man enter into the temple: that is, could penetrate the secret, until the seven plagues of the seven angels were finished. Wherefore, the Psalmist: “This,” he says, “is labor before me, until I may enter into the sanctuary of God and understand the conclusion of matters, and so forth.” Here is signified therefore, that saints before the judgment will not know the secret mysteries of God’s judgments. Let it then suffice us that he himself has deigned to reveal to us; for the rest, let us believe that the Lord is just in all his ways and holy in all his works. To him be glory.
